\begin{abstract}
Mathematical reasoning remains a formidable challenge for language models, given its inherently intricate and highly structured characteristics. In this study, we present DeepSeekMath~7B, developed by further pre-training DeepSeek-Coder-Base-v1.5 7B with 120B mathematics-focused tokens obtained from Common Crawl, in addition to natural language and code corpora. DeepSeekMath~7B attains a remarkable 51.7\% accuracy on the competition-level MATH benchmark without the aid of external toolkits or ensemble strategies, nearing the proficiency of proprietary models such as Gemini-Ultra and GPT-4. When applying self-consistency across 64 responses, DeepSeekMath~7B’s performance increases to 60.9\% on MATH. The model's enhanced mathematical reasoning proficiency can be credited to two principal innovations: firstly, the construction of an elaborate data curation pipeline leveraging extensive open-access web resources; secondly, the development of Group Relative Policy Optimization (GRPO)—an adaptation of Proximal Policy Optimization (PPO)—which boosts mathematical reasoning capabilities while simultaneously improving PPO's memory efficiency.
\end{abstract}

\section{Introduction}

Large language models (LLMs) have dramatically transformed artificial intelligence methods for mathematical reasoning, yielding significant progress in quantitative reasoning benchmarks \citep{MATH} as well as geometry-focused benchmarks \citep{trinh2024solving}. Furthermore, these models have demonstrated substantial utility in supporting humans with challenging mathematical problem-solving tasks \citep{tao}. Nonetheless, state-of-the-art systems like GPT-4 \citep{gpt4} and Gemini-Ultra \citep{gemini} remain closed-source, leaving available open-source alternatives notably behind in terms of performance.

This paper presents DeepSeekMath, a specialized language model for mathematics that markedly surpasses the abilities of other open-source systems and nearly matches the performance of GPT-4 on standard academic tests. To realize this advance, we constructed the DeepSeekMath Corpus—a large, high-quality pre-training dataset containing 120B mathematical tokens. The corpus is sourced from Common Crawl (CC) with the help of a fastText-based classifier \citep{joulin2016fasttext}. Initially, the classifier is built by training on positive samples from OpenWebMath \citep{openwebmath} and using a diverse array of unrelated web content as negative samples. In subsequent iterations, the classifier identifies further relevant examples in the CC, which are then vetted and refined through manual annotation. We then retrain the classifier with this augmented dataset to enhance its selection accuracy. Our evaluations reveal that the resultant large-scale dataset maintains a high standard, as demonstrated by DeepSeekMath-Base 7B achieving scores of 64.2\% on GSM8K \citep{gsm8k} and 36.2\% on the MATH competition dataset \citep{MATH}, both figures surpassing those reported for Minerva 540B \citep{minerva}. Furthermore, since the DeepSeekMath Corpus includes multilingual content, we observe performance gains on Chinese mathematics benchmarks \citep{wei2023cmath,agieval}. We posit that our methodology for assembling mathematical training data offers a useful foundation for ongoing work in this field, with ample opportunities for future refinement and innovation.

DeepSeekMath-Base leverages DeepSeek-Coder-Base-v1.5 7B \citep{deepseek-coder} as its foundation, based on the observation that initiating training from a code-centric model yields superior results compared to general-purpose LLMs. Additionally, our findings suggest that mathematical training not only elevates mathematical proficiency but also strengthens general reasoning abilities, as evidenced by performance gains on the MMLU \citep{mmlu} and BBH \citep{bbh} benchmarks.

Following pre-training, DeepSeekMath-Base undergoes specialized mathematical instruction tuning, utilizing datasets for chain-of-thought \citep{cot}, program-of-thought \citep{pot,pal}, and tool-augmented reasoning \citep{tora}. The outcome, DeepSeekMath-Instruct 7B, surpasses all other 7B-scale models and reaches a level of performance on par with leading 70B instruction-tuned open-source models.

Moreover, we present Group Relative Policy Optimization (GRPO), a novel variant of reinforcement learning (RL) rooted in Proximal Policy Optimization (PPO) \citep{schulman2017proximal}. Unlike PPO, GRPO omits the use of a critic, instead employing group-based score baselines to markedly cut down on required computational resources. When trained solely on selected English instruction tuning data, GRPO markedly enhances the already robust DeepSeekMath-Instruct, producing notable gains both in domain (GSM8K: 82.9\% $\rightarrow$ 88.2\%, MATH: 46.8\% $\rightarrow$ 51.7\%) and across out-of-domain mathematical benchmarks (e.g., CMATH: 84.6\% $\rightarrow$ 88.8\%) during RL optimization. Additionally, we establish a unified framework that integrates various approaches, including Rejection Sampling Fine-Tuning (RFT) \citep{yuan2023scaling}, Direct Preference Optimization (DPO) \citep{dpo}, PPO, and GRPO. Within this unified perspective, we observe that these approaches can be interpreted as direct or simplified variants of RL algorithms. To further examine this framework, we perform thorough experiments, such as comparing online versus offline learning, outcome versus process supervision, and single-turn against iterative RL approaches, thereby elucidating critical components of this paradigm. Finally, we articulate the mechanisms underlying the performance improvements in instruction-tuned models due to RL, and discuss promising avenues for advancing RL effectiveness within this unified schema.

\subsection{Contributions}
This work presents advancements in large-scale mathematical pre-training and investigates reinforcement learning approaches.

\textbf{Math Pre-Training at Scale}

\begin{itemize}[topsep=0pt]
    \item We demonstrate that Common Crawl, a widely available web dataset, possesses substantial mathematical content. Utilizing a carefully crafted data selection pipeline, we curate the DeepSeekMath Corpus, a high-fidelity dataset comprising 120B tokens sourced from math-focused web pages. This corpus is nearly 7 times larger than that employed by Minerva \citep{minerva} and about 9 times greater than the recent OpenWebMath release \citep{openwebmath}.

    \item The DeepSeekMath-Base 7B, our foundational pre-trained model, delivers performance on par with the much larger Minerva 540B \citep{minerva}. This result suggests that parameter count alone does not dictate mathematical reasoning effectiveness; instead, compact models, when pre-trained on high-quality datasets, can excel.

    \item Our experimental results in math pre-training indicate that models benefit significantly from being trained on code before exposure to mathematical data, enhancing their capability to resolve math problems both with and without auxiliary tools. This provides partial insight into the ongoing debate on whether code pre-training augments reasoning—at least in mathematical contexts, it appears beneficial.

    \item We observe that, despite the prevalence of arXiv papers in mathematical data pipelines, incorporating arXiv into pre-training does not yield marked gains on any mathematical benchmark considered herein.
\end{itemize}

\textbf{Exploration and Analysis of Reinforcement Learning}

\begin{itemize}[topsep=0pt]
    \item This work presents Group Relative Policy Optimization (GRPO), a novel reinforcement learning approach designed for both efficacy and efficiency. Unlike conventional methods that utilize a critic network, GRPO leverages group score-based baseline estimation, enabling substantial savings in computational resources relative to Proximal Policy Optimization (PPO).
    \item Our experiments reveal that GRPO substantially improves the results of the instruction-tuned model DeepSeekMath-Instruct using solely instruction-tuning data. Additionally, we find that reinforcement learning via GRPO leads to marked gains in out-of-domain generalization.
    \item We establish a cohesive framework that unifies various techniques—such as RFT, DPO, PPO, and GRPO—providing conceptual clarity among them. Moreover, we systematically examine factors such as online versus offline training, supervision at the outcome versus process level, and single-turn versus iterative reinforcement learning, thus uncovering key ingredients of the overall framework through rigorous experimentation.
    \item Utilizing our unified framework, we investigate underlying mechanisms responsible for the empirical strengths of reinforcement learning, and outline several promising avenues to further enhance LLMs through improved reinforcement learning techniques.
\end{itemize}

\subsection{Summary of Evaluations and Metrics}

\begin{itemize}[topsep=0pt]
    \item \textbf{English and Chinese Mathematical Reasoning}: Our models undergo a thorough evaluation on a spectrum of English and Chinese benchmarks that span mathematical reasoning tasks from elementary to collegiate levels. The English testbeds comprise GSM8K \citep{gsm8k}, MATH \citep{MATH}, SAT \citep{llemma}, OCW Courses \citep{minerva}, and MMLU-STEM \citep{mmlu}; the Chinese suite includes MGSM-zh \citep{mgsm}, CMATH \citep{wei2023cmath}, Gaokao-MathCloze \citep{agieval}, and Gaokao-MathQA \citep{agieval}. We measure model proficiency both in crafting comprehensive natural language solutions without external tools and in resolving problems via Python code execution.

    For English datasets, DeepSeekMath-Base demonstrates competitive performance relative to the proprietary Minerva 540B \citep{minerva}, and consistently outperforms all open-source foundational models—including Mistral 7B \citep{mistral} and Llemma-34B \citep{llemma}—regardless of their math-oriented pre-training backgrounds, frequently by considerable margins. The advantage of DeepSeekMath-Base is even more pronounced for Chinese datasets, which we attribute to our use of high-quality multilingual mathematical pre-training data in contrast to the English-only focus adopted by prior work \citep{minerva,llemma}. By applying both mathematical instruction tuning and reinforcement learning, DeepSeekMath-Instruct and DeepSeekMath-RL achieve state-of-the-art results, breaking the 50\% accuracy threshold on the challenging competition-grade MATH benchmark—a first among open-source efforts.
\end{itemize}

\item \textbf{Formal Mathematics}: We assess the capabilities of DeepSeekMath-Base in the context of informal-to-formal theorem proving, following the methodology outlined by \citep{dsp_proof}, utilizing the miniF2F benchmark \citep{minif2f} with Isabelle \citep{isabelle} as the proof assistant. The results indicate that DeepSeekMath-Base achieves robust few-shot performance in autoformalization tasks.
\item \textbf{Natural Language Understanding, Reasoning, and Code}: To obtain a holistic evaluation of the model’s linguistic comprehension, reasoning, and programming proficiency, DeepSeekMath-Base is tested on the Massive Multitask Language Understanding (MMLU) benchmark \citep{mmlu}, encompassing 57 diverse multiple-choice problems, as well as BIG-Bench Hard (BBH) \citep{bbh}, which consists of 23 problems emphasizing multi-step reasoning, and coding tasks including HumanEval \citep{codex} and MBPP \citep{mbpp}, widely adopted for code model assessment. Results suggest that math-oriented pre-training leads to improvements in both general language understanding and reasoning accuracy.
\end{itemize}

\section{Math Pre-Training}

\subsection{Data Collection and Decontamination} This section details the methodology for assembling the DeepSeekMath Corpus from Common Crawl. As illustrated in Figure \ref{fig:data_collect}, the process employs an iterative pipeline, systematically curating a large-scale mathematics dataset starting from a limited but high-quality initial seed corpus, such as a specialized collection of mathematical data. This method can be readily adapted for other fields, including code datasets.

The process begins by designating OpenWebMath \citep{openwebmath}—a dataset of rigorously curated mathematical web materials—as the seed. A fastText model \citep{joulin2016fasttext} is trained to identify web pages with similar mathematical characteristics. To construct the training set, 500,000 entries from the seed corpus are chosen as positive examples and an additional 500,000 Common Crawl web pages are sampled as negative examples. Training is conducted using an open-source fastText implementation\footnote{\url{https://fasttext.cc}}, with the following hyperparameters: a vector size of 256, a learning rate of 0.1, maximum word n-gram length of 3, a minimum occurrence threshold of 3 words, and 3 epochs. Prior to recall, the scale of Common Crawl is reduced through URL-based deduplication and near-duplicate filtering, yielding a corpus of 40B HTML documents. Candidate mathematical webpages are retrieved using the fastText model and are subsequently ranked by model-assigned score; only those receiving the highest scores are retained. The final volume to preserve is empirically determined via pre-training studies, considering token subsets of the top 40B, 80B, 120B, and 160B tokens. In the initial round, the top 40B tokens are selected for further use.

Following the initial round of data acquisition, a significant number of mathematical web pages remain unretrieved. This shortfall is primarily due to the limited diversity of positive samples in the initial fastText model’s training data. To address this, we augment the seed corpus by sourcing more mathematical web domains, thereby enhancing the fastText model’s capacity. In detail, we segment the entire Common Crawl dataset into non-overlapping domains, where each domain encompasses web pages sharing an identical base URL. For every domain, we compute the proportion of web pages collected during the first pass. Domains with a collection rate exceeding 10\% are designated as math-oriented (such as \url{mathoverflow.net}).

Next, we undertake a manual review to label URLs within these domains that correspond to mathematical resources (e.g., \url{mathoverflow.net/questions}). Pages linked to these mathematical URLs that were missed previously are then incorporated into the seed corpus. Through this procedure, we amass additional positive training examples, enabling us to fine-tune the fastText model for greater mathematical content recall in the subsequent cycle. After four rounds, the process results in a collection of 35.5M mathematical web pages, amassing 120B tokens in total. By the fourth iteration, approximately 98\% of the data has already been acquired by the third iteration, leading us to halt further data collection at this point.

To prevent any leakage from benchmarks, we adopt the method from \cite{deepseek-coder} to exclude web pages containing either questions or answers sourced from established English mathematical benchmarks such as GSM8K \citep{gsm8k} and MATH \citep{MATH}, as well as Chinese datasets such as CMATH \citep{wei2023cmath} and AGIEval \citep{agieval}. The filtering process entails removing any text where a sequence of 10 consecutive tokens matches a substring from any of the aforementioned benchmarks. In cases where benchmark text is shorter than 10 tokens but includes at least 3, we apply exact string matching to eliminate possibly contaminated web content from the math training set.

\subsection{Validating the Quality of the DeepSeekMath~Corpus}

We conduct pre-training assessments to evaluate the performance of the DeepSeekMath~Corpus relative to other recently introduced math-oriented corpora:

\begin{itemize}[topsep=0pt]
    \item \textbf{MathPile} \citep{mathpile}: a composite corpus comprising 8.9B tokens sourced from textbooks, Wikipedia, ProofWiki, CommonCrawl, StackExchange, and arXiv, with arXiv accounting for over 85\% of the dataset;
    \item \textbf{OpenWebMath} \citep{openwebmath}: a dataset formed by extracting math-relevant material from CommonCrawl, consisting of 13.6B tokens;
    \item \textbf{Proof-Pile-2} \citep{llemma}: this collection merges OpenWebMath, AlgebraicStack (10.3B tokens of mathematical code), and arXiv manuscripts (28.0B tokens). In our experiments with Proof-Pile-2, we adhere to the arXiv:Web:Code composition ratio of 2:4:1 as stipulated in \cite{llemma}.
\end{itemize}

\subsubsection{Training Setting}
\label{sec:quality-policy}
We perform mathematical fine-tuning on a 1.3B-parameter general-purpose pre-trained language model, utilizing the architecture of DeepSeek LLMs \citep{deepseek-llm}; this variant is referenced as DeepSeek-LLM 1.3B. For each mathematical dataset, we train the model independently for 150B tokens. The training process is executed via the resource-efficient HAI-LLM framework \citep{haillm}. Adopting the methodology established for DeepSeek LLMs, optimization relies on AdamW \citep{adamW} with $\beta_1=0.9$, $\beta_2=0.95$, and a $\mathrm{weight\_decay}$ of $0.1$. A multi-phase learning rate strategy is applied: the learning rate peaks after an initial 2,000-step warmup, subsequently diminishes to $31.6\%$ at $80\%$ of the total steps, and declines further to $10.0\%$ of the peak by $90\%$ of training. The peak learning rate is set at $5.3\times 10^{-4}$, employing a batch size encompassing 4M tokens and a 4K token context window.

\subsubsection{Evaluation Results}

\textbf{The DeepSeekMath~Corpus distinguishes itself as the largest collection, providing extensive multilingual mathematical data of notably high quality.}

\begin{itemize}[topsep=0pt]
    \item \textbf{High-quality}:
    We assess the model's ability on 8 mathematical evaluation datasets using few-shot chain-of-thought prompting \cite{cot}. According to Table \ref{tab:corpora_comparison}, the model trained with DeepSeekMath~Corpus exhibits a distinct superiority in performance. Furthermore, as illustrated in Figure \ref{fig:corpora_comparisons}, models trained on DeepSeekMath~Corpus consistently outperform those trained on Proof-Pile-2 at the 50B token mark (representing one full epoch of Proof-Pile-2), thereby affirming the superior average quality of the DeepSeekMath~Corpus.
    \item \textbf{Multilingual}:
    DeepSeekMath~Corpus incorporates content in a variety of languages, with English and Chinese as the most frequent. Results in Table \ref{tab:corpora_comparison} demonstrate that models trained on DeepSeekMath~Corpus yield stronger mathematical reasoning abilities in both English and Chinese. By comparison, other mathematical datasets, which predominantly focus on English, offer minimal gains or potentially detract from Chinese mathematical reasoning performance.
    \item \textbf{Large-scale}:
    DeepSeekMath~Corpus substantially exceeds existing mathematical datasets in terms of size. Figure \ref{fig:corpora_comparisons} illustrates that DeepSeek-LLM 1.3B, after training on DeepSeekMath~Corpus, exhibits a sharper and more sustained progression in learning. Baseline corpora, being much more limited in scale and reused several times during training, result in models whose performance plateaus considerably sooner.
\end{itemize}

\subsection{Training and Evaluating DeepSeekMath-Base 7B}

This section presents DeepSeekMath-Base 7B, a foundational model tailored for robust mathematical reasoning. The model is built upon DeepSeek-Coder-Base-v1.5 7B \citep{deepseek-coder}, serving as its initialization, and subsequently trained on a total of 500B tokens. The composition of the training dataset is as follows: 56\% derived from the DeepSeekMath Corpus, 4\% sourced from AlgebraicStack, 10\% from arXiv, 20\% from Github code repositories, and the final 10\% comprising natural language text in both English and Chinese from Common Crawl. While the primary training regimen follows the guidelines described in Section \ref{sec:quality-policy}, modifications include raising the peak learning rate to 4.2e-4 and utilizing a batch size of 10M tokens.

To rigorously evaluate the mathematical proficiency of DeepSeekMath-Base 7B, we analyze its capacity to autonomously generate comprehensive mathematical solutions, leverage tools to address mathematical tasks, and perform formal theorem verification. Additionally, the model’s broader competencies—such as its understanding of natural language, logical reasoning, and programming proficiency—are also systematically examined.

\paragraph{Step-by-Step Mathematical Problem Solving}

For benchmarking the stepwise problem-solving aptitude of DeepSeekMath-Base, we employ few-shot chain-of-thought prompting \citep{cot} on a set of eight diverse datasets in both English and Chinese. These benchmarks are drawn from various mathematical disciplines and include quantitative reasoning tasks (e.g., GSM8K \citep{gsm8k}, MATH \citep{MATH}, CMATH \citep{wei2023cmath}) as well as multiple-choice assessments (e.g., MMLU-STEM \citep{mmlu} and Gaokao-MathQA \citep{agieval}), collectively spanning elementary through college-level mathematical challenges.

As indicated in Table \ref{tab:base_math_cot}, DeepSeekMath-Base 7B achieves the highest results among open-source base models on all eight evaluated benchmarks. This comparison includes the widely adopted Mistral 7B \citep{mistral} and the more recent Llemma 34B \citep{llemma}, the latter having been trained on mathematical data from Proof-Pile-2 \citep{llemma}. Particularly remarkable is its performance on the MATH benchmark, representative of competition-level problem-solving, where DeepSeekMath-Base surpasses other open-source base models by more than 10\% in absolute terms, and even exceeds Minerva 540B \citep{minerva}—a proprietary model of substantially greater scale that builds upon PaLM \citep{palm} and is fine-tuned on mathematical corpora.

\paragraph{Mathematical Problem Solving with Tool Use} To assess the capability for program-assisted mathematical reasoning, we conduct evaluations on GSM8K and MATH using few-shot program-of-thought prompting \citep{pot,pal}. Each task prompts models to generate a Python script for the problem at hand, with access to packages such as \textit{math} and \textit{sympy} to handle complex operations. The answer is determined based on the output of these scripts. Table \ref{tab:base_math_pot_proof} demonstrates that DeepSeekMath-Base 7B sets a new state-of-the-art among open-source models, outperforming the previous leading model, Llemma 34B.

\paragraph{Formal Mathematics}

Automating formal proof construction enhances both the reliability and efficiency of mathematical reasoning, an area of rising research interest. We assess DeepSeekMath-Base 7B on the informal-to-formal proving task from \citep{dsp_proof}, wherein the model generates formal proofs from informal statements and corresponding formal/informal proofs. Evaluations are performed on the miniF2F dataset \citep{minif2f}, which targets formal Olympiad-level math; here, the model is prompted with few-shot examples to produce formal proofs in Isabelle for each item. Following the protocol in \cite{dsp_proof}, proof sketches generated by the model are supplemented using Sledgehammer \citep{sledgehammer}, an automated proving tool that completes missing steps. Table \ref{tab:base_math_pot_proof} shows that DeepSeekMath-Base 7B achieves notable results in automatic formalization of proofs.

\paragraph{Natural Language Understanding, Reasoning, and Code}

For natural language comprehension, reasoning, and programming, model abilities are measured using MMLU \citep{mmlu}, BBH \citep{bbh}, and coding-focused evaluations HumanEval \citep{codex} and MBPP \citep{mbpp}. As illustrated in Table \ref{tab:understanding_reasoning_code}, DeepSeekMath-Base 7B delivers substantial gains over its predecessor, DeepSeek-Coder-Base-v1.5 \citep{deepseek-coder}, in MMLU and BBH, attesting to the positive influence of mathematical training on linguistic and reasoning proficiencies. Furthermore, through the continued incorporation of code tokens during training, DeepSeekMath-Base 7B preserves the strong code generation abilities of DeepSeek-Coder-Base-v1.5 on HumanEval and MBPP. On all three benchmarks evaluating reasoning and coding, DeepSeekMath-Base 7B consistently outperforms the generalist Mistral 7B \citep{mistral}.

\section{Supervised Fine-Tuning}

\subsection{SFT Data Curation}

We develop a comprehensive instruction-tuning dataset for mathematics that includes problems written in both English and Chinese, spanning multiple mathematical disciplines and varying degrees of complexity. Each problem is aligned with solutions in the chain-of-thought (CoT) format \citep{cot}, the program-of-thought (PoT) format \citep{pot,pal}, and a tool-integrated reasoning format \citep{tora}. The dataset comprises 776K training samples in total.

\begin{itemize}[topsep=0pt]
    \item \textbf{English mathematical datasets}: Our annotation efforts extend to providing tool-integrated solutions for GSM8K and MATH problems. We also incorporate a selected portion of MathInstruct \citep{MathInstruct} and use the Lila-OOD \citep{lila} training data, in which solutions are given through either CoT or PoT. This portion of the corpus encompasses a broad spectrum of mathematical fields, such as algebra, probability, number theory, calculus, and geometry.
    \item \textbf{Chinese mathematical datasets}: For Chinese, we assemble a collection of K-12 mathematics questions from 76 distinct sub-domains including, for example, linear equations. Each question is supplemented with solutions written in both CoT and tool-integrated reasoning styles.
\end{itemize}

\subsection{Training and Evaluating DeepSeekMath-Instruct 7B}

This section presents DeepSeekMath-Instruct 7B, a model that receives instruction fine-tuning based on DeepSeekMath-Base. For each batch, training samples are concatenated at random up to a maximum context window of 4K tokens. The fine-tuning procedure involves 500 steps, utilizing a batch size of 256 and maintaining a fixed learning rate of 5e-5 throughout.

Model performance in mathematics is assessed with and without tool assistance, leveraging four quantitative reasoning benchmarks in both English and Chinese. To contextualize our results, we compare our model's capabilities against the prevailing state-of-the-art systems:

\begin{itemize}[topsep=0pt]
    \item \textbf{Closed-source models}: This group includes the GPT series, specifically GPT-4 \citep{gpt4} and GPT-4 Code Interpreter~\footnote{\url{https://openai.com/blog/chatgpt-plugins\#\#code-interpreter}}, as well as Gemini Ultra and Pro \citep{gemini}, Inflection-2 \citep{inflection-2}, Grok-1~\footnote{\url{https://x.ai/model-card}}, and recent Chinese releases such as Baichuan-3~\footnote{\url{https://www.baichuan-ai.com}} and the most recent GLM-4~\footnote{\url{https://open.bigmodel.cn/dev/api\#glm-4}} from the GLM series \citep{glm}.
\end{itemize}

These models serve general-purpose applications and typically undergo extensive alignment processes.
    \item \textbf{Open-source models} comprise a range of systems, such as: (1) DeepSeek-LLM-Chat 67B \citep{deepseek-llm}, (2) Qwen 72B \citep{qwen}, (3) SeaLLM-v2 7B \citep{seallm}, and (4) ChatGLM3 6B \citep{chatglm3}, which are broadly applicable, alongside mathematics-focused variants, including (5) InternLM2-Math 20B~\footnote{\url{https://github.com/InternLM/InternLM-Math}}, an extension of InternLM2 subject to mathematical training and subsequent instruction fine-tuning, (6) Math-Shepherd-Mistral 7B, developed through PPO-based optimization \citep{schulman2017proximal} of Mistral 7B \citep{mistral} leveraging a process-supervised reward signal, (7) the WizardMath collection \citep{wizardmath}, which augments mathematical capabilities in Mistral 7B and Llama-2 70B \citep{llama2} via evolve-instruct (a variant of instruction tuning with AI-curated tasks) and PPO reinforcement, primarily using GSM8K and MATH datasets for training, (8) MetaMath 70B \citep{metamath}, a Llama-2 70B checkpoint further fine-tuned on enriched GSM8K and MATH data, (9) ToRA 34B \cite{tora}, produced by further training CodeLlama 34B for tool-aided mathematical reasoning, and (10) MAmmoTH 70B \citep{MathInstruct}, where Llama-2 70B is instruction-tuned utilizing MathInstruct data.
\end{itemize}

According to the results summarized in Table \ref{tab:sft_rl_math}, when tool-assisted approaches are not permitted, DeepSeekMath-Instruct 7B achieves robust stepwise reasoning capabilities. Specifically, on the challenging MATH benchmark, this model exceeds the performance of all other open-source alternatives and outperforms many proprietary systems (such as Inflection-2 and Gemini Pro) by a margin of at least 9\% in absolute terms. This remains true even relative to substantially larger models (for instance, Qwen 72B) or those trained with specialized mathematical RL techniques (e.g., WizardMath-v1.1 7B). While DeepSeekMath-Instruct is comparable to proprietary Chinese models such as GLM-4 and Baichuan-3 on MATH, it does not reach the proficiency of GPT-4 or Gemini Ultra.

In a setting where models can integrate both natural language reasoning and programmatic tool usage to solve problems, DeepSeekMath-Instruct 7B attains nearly 60\% accuracy on MATH—surpassing all existing open-source approaches. On additional evaluation datasets, the model performs on par with DeepSeek-LLM-Chat 67B, the prior state-of-the-art model, despite being only one-tenth the size.

\section{Reinforcement Learning}

\subsection{Group Relative Policy Optimization} Reinforcement learning (RL) has proven to further boost the mathematical problem-solving capabilities of large language models after supervised fine-tuning (SFT) \citep{wang2023math,wizardmath}. In this section, we present Group Relative Policy Optimization (GRPO), our RL framework designed for both efficacy and efficiency.

\subsubsection{From PPO to GRPO} Proximal Policy Optimization (PPO) \citep{schulman2017proximal} stands out as a commonly adopted actor-critic RL algorithm in the RL fine-tuning phase of LLMs \citep{ouyang2022training}. Specifically, PPO enhances model parameters by maximizing the surrogate objective given as follows:

\begin{equation}
\footnotesize
    \mathcal{J}_{PPO}(\theta) = \mathbb{E}{[q \sim P(Q), o \sim \pi_{\theta_{old}}(O|q)]} \frac{1}{|o|} \sum_{t=1}^{|o|} \min \left[ \frac{\pi_\theta(o_{t} | q, o_{<t})}{\pi_{\theta_{old}}(o_{t} | q, o_{<t})} A_{t}, \text{clip} \left( \frac{\pi_\theta(o_{t} | q, o_{<t})}{\pi_{\theta_{old}}(o_{t} | q, o_{<t})

Here, $\pi_{\theta}$ and $\pi_{\theta_{old}}$ denote the current and preceding policy models, respectively. The variables $q$ and $o$ represent questions and outputs that are sampled from the question corpus and generated according to the old policy $\pi_{\theta_{old}}$. The parameter $\varepsilon$ serves as a clipping hyper-parameter utilized in PPO to stabilize the learning process. The term $A_t$ refers to the advantage, computed using Generalized Advantage Estimation (GAE) \citep{gae}, which leverages the sequence of rewards $\{r_{\ge t}\}$ in conjunction with an estimated value function $V_{\psi}$. As a result, PPO requires the concurrent training of a value function along with the policy network. To prevent excessive optimization towards the reward model, the conventional method includes a per-token Kullback-Leibler (KL) divergence penalty, relative to a fixed reference policy, incorporated directly into the reward for each token \citep{ouyang2022training}, formulated as

\begin{equation}
    r_{t} = r_\varphi(q, o_{\le t}) - \beta \log\frac{\pi_{\theta}(o_{t}|q, o_{<t})}{\pi_{ref}(o_{t}|q, o_{<t})},
\label{eq:PPO-reward}
\end{equation}

where $r_\varphi$ stands for the reward model, $\pi_{ref}$ refers to the reference policy (typically, the initial SFT model), and $\beta$ is the scaling factor for the KL term.

Given that the value function in PPO generally constitutes a model of similar scale as the main policy, this introduces notable memory usage and computational overhead. Furthermore, during reinforcement learning, the value function operates as a baseline in advantage estimation to help reduce variance. However, in large language model settings, it is standard for only the final token to be assigned a reward value by the reward model, making accurate value function estimation for each token non-trivial. To address these limitations, as depicted in Figure \ref{fig:grpo}, we introduce Group Relative Policy Optimization (GRPO). GRPO eliminates the need for auxiliary value function approximation present in PPO, substituting it with the mean reward across a set of sampled outputs—generated in response to the same prompt—as the baseline. Specifically, for each question $q$, GRPO samples a collection of outputs $\{o_1, o_2, \cdots, o_G\}$ from $\pi_{\theta_{old}}$, and then seeks to optimize the policy by maximizing the following objective:

\begin{equation}
\footnotesize
\begin{split}
    \mathcal{J}_{GRPO}(\theta) &= \mathbb{E}{[q \sim P(Q), \{o_i\}_{i=1}^G \sim \pi_{\theta_{old}}(O|q)]}  \\
    & \frac{1}{G}\sum_{i=1}^G\frac{1}{|o_i|} \sum_{t=1}^{|o_i|} \left\{ \min \left[ \frac{\pi_\theta(o_{i,t} | q, o_{i,<t})}{\pi_{\theta_{old}}(o_{i,t} | q, o_{i,<t})} \hat{A}_{i,t}, \text{clip} \left( \frac{\pi_\theta(o_{i,t} | q, o_{i,<t})}{\pi_{\theta_{old}}(o_{i,t} | q, o_{i,<t})}, 1 - \varepsilon, 1 + \varepsilon \right)  \hat{A}_{i,t} \right] - \beta \mathbb{D}_{KL}\left[\pi_{\theta} || \pi_{ref}\right]\right\} ,
\end{split}
\label{eq:GRPO-obj}
\end{equation}

In this expression, $\varepsilon$ and $\beta$ are hyper-parameters, and $\hat{A}_{i,t}$ denotes the advantage value computed exclusively based on the relative performance among the group’s outputs—a procedure which is further explained in the subsequent sections. The design of GRPO for calculating advantages based on output groups aligns naturally with the comparative foundation of reward models, which are typically trained on datasets featuring pairwise comparisons for the same prompt. It is also notable that, instead of integrating the KL term directly into the reward, GRPO regularizes training by explicitly including the KL divergence between the trained policy and the reference policy as a loss penalty, thereby simplifying the computation of $\hat{A}_{i,t}$. Moreover, in contrast to

\begin{algorithm}[t]
  \small
  \caption{Iterative Group Relative Policy Optimization}
  \textbf{Input} initial policy model $\pi_{\theta_{\text{init}}}$; reward models $r_\varphi$; task prompts $\mathcal{D}$; 
  hyperparameters $\varepsilon$, $\beta$, $\mu$
  \begin{algorithmic}[1]
    \State Set policy model $\pi_\theta \leftarrow \pi_{\theta_{\text{init}}}$
    \For{iteration = 1, \dots, I}
       \State Assign reference model $\pi_{ref} \leftarrow \pi_{\theta}$
      \For{step = 1, \dots, M}
      \State Draw a batch $\mathcal{D}_b$ from $\mathcal{D}$
      \State Update old policy model $\pi_{\theta_{old}} \leftarrow \pi_{\theta}$ 
      \State For each question $q \in \mathcal{D}_b$, generate $G$ responses $\{o_i\}_{i=1}^G \sim \pi_{\theta_{old}} (\cdot \mid q) $
      \State Evaluate rewards $\{r_i\}_{i=1}^{G}$ for generated outputs $o_i$ using $r_{\varphi}$
      \State Determine $\hat{A}_{i,t}$ for the $t$-th token in $o_i$ utilizing group relative advantage estimation
      \For{GRPO iteration = 1, \dots, $\mu$}
        \State Update policy model $\pi_{\theta}$ to maximize the GRPO objective in Equation \ref{eq:GC-GRPO}
      \EndFor
    \EndFor 
    \State Perform continual training of $r_\varphi$ via a replay-based update. 
    \EndFor 
  \end{algorithmic}
  \textbf{Output} $\pi_\theta$
  \label{alg:iter-grpo}
\end{algorithm}

\subsubsection{Outcome Supervision RL with GRPO} In this approach, for a given question $q$, a set of candidate outputs $\{o_1, o_2, \cdots, o_G\}$ are produced from the previous policy $\pi_{\theta_{old}}$. Each output is evaluated by the reward model, yielding a reward vector $\mathbf{r} = \{r_1, r_2, \cdots, r_G\}$. To normalize, each reward is centered by subtracting the group mean and then scaled by dividing by the group standard deviation. Under outcome supervision, the reward is assigned in normalized form only at the final token of each output $o_i$, and every token $t$ within the output inherits this normalized terminal reward as its advantage estimate, i.e., $\hat{A}_{i, t} = \widetilde{r}_i = \frac{r_i- {\rm mean}(\mathbf{r})}{{\rm std}(\mathbf{r})}$. Policy parameters are then updated via maximization of the objective in equation (\ref{eq:GRPO-obj}).

\subsubsection{Process Supervision RL with GRPO} Outcome supervision restricts feedback to the sequence’s end, which can be insufficient for directing learning in intricate mathematical reasoning scenarios. Building on \cite{wang2023math}, process supervision extends the reward assignment to every step in the reasoning chain. Specifically, given a question $q$ and group of $G$ outputs $\{o_1, o_2, \cdots, o_G\}$, a process reward model assigns a score to each step, resulting in the reward set: $\mathbf{R} = \{ \{r_1^{index(1)},\cdots,r_1^{index(K_1)}\}, \cdots,  \{r_G^{index(1)},\cdots,r_G^{index(K_G)}\} \}$, where $index(j)$ specifies the final token position of the $j$-th step, and $K_i$ denotes the count of steps for output $o_i$. Rewards are normalized analogously: $\widetilde{r}_i^{index(j)} = \frac{r_i^{index(j)} - {\rm mean(\mathbf{R})}}{{\rm std(\mathbf{R})}}$.
Next, for each token, process supervision assigns the advantage as the cumulative sum of normalized future step rewards: $\hat{A}_{i, t} = \sum_{index(j) \ge t} \widetilde{r}_i^{index(j)}$.
Policy optimization proceeds by maximizing the same objective as in equation (\ref{eq:GR

\subsubsection{Iterative RL with GRPO} As reinforcement learning proceeds, earlier versions of the reward model may become inadequate for effectively supervising the evolving policy model. To address this, we implement an iterative RL strategy employing GRPO. As illustrated in Algorithm \ref{alg:iter-grpo}, the iterative GRPO approach involves creating updated datasets for the reward model using samples generated by the current policy model. The reward model undergoes continuous training utilizing a replay buffer that preserves 10\% of previous data. Subsequently, we replace the reference model with the updated policy model and proceed to further train the policy model with the revised reward model.

\subsection{Training and Evaluating DeepSeekMath-RL}

Our RL experiments are built upon DeepSeekMath-Instruct 7B. The RL training data comprises approximately 144K chain-of-thought-format questions sourced from GSM8K and MATH within the SFT dataset. All other SFT questions are omitted to isolate and measure the effects of RL on benchmarks not represented during the RL training process. Following the methodology described in \citep{wang2023math}, we construct the dataset for reward model training. The initial reward model is trained from DeepSeekMath-Base 7B using a learning rate of 2e-5. In the GRPO procedure, the policy model is trained with a learning rate set to 1e-6, and the KL regularization coefficient is fixed at 0.04. We draw $64$ candidate answers per question. The maximum sequence length is configured at 1024 tokens, with a batch size also set to 1024. Each exploration phase is followed by a single policy model update. For evaluation, DeepSeekMath-RL 7B is assessed on the same benchmarks as DeepSeekMath-Instruct 7B. In this context, chain-of-thought GSM8K and MATH tasks are categorized as in-domain, whereas all other benchmarks are classified as out-of-domain.

Table \ref{tab:sft_rl_math} presents a comparative analysis of open- and closed-source models using both chain-of-thought and tool-augmented reasoning strategies on benchmarks in English and Chinese. The following observations emerge: 1) DeepSeekMath-RL 7B achieves accuracies of 88.2\% on GSM8K and 51.7\% on MATH through chain-of-thought reasoning. These results not only exceed the performance of all open-source models within the 7B to 70B range, but also outpace most closed-source alternatives. 2) Importantly, DeepSeekMath-RL 7B was exclusively trained on chain-of-thought instruction tuning data drawn from GSM8K and MATH, with initialization from DeepSeekMath-Instruct 7B. Despite this narrow training focus, it surpasses DeepSeekMath-Instruct 7B on every evaluated metric, underscoring the significant impact of reinforcement learning.

\section{Discussion}
In the following section, we detail the insights gleaned from our pre-training and reinforcement learning studies.

\subsection{Lessons Learnt in Pre-Training}

We begin by discussing key takeaways from pre-training. Unless stated otherwise, our discussion assumes the training procedures specified in Section \ref{sec:quality-policy}. Specifically, references to the DeepSeekMath Corpus pertain to an 89B-token dataset from the second round of data collection.

\subsubsection{Code Training Benefits Mathematical Reasoning}

The hypothesis that training on code enhances reasoning ability is widespread but has yet to be definitively proven. Here, we seek to partially address this question, focusing on mathematical tasks: code training enhances a model's capability for mathematical reasoning, whether or not auxiliary tools are employed.

To systematically explore the influence of code pre-training on mathematical reasoning, we designed experiments comparing two approaches: a sequential two-stage training process versus a one-stage training regime.

\textbf{Two-Stage Training}

\begin{itemize}[topsep=0pt]
    \item \textbf{Code Training for 400B Tokens $\rightarrow$ Math Training for 150B Tokens}: DeepSeek-LLM 1.3B is first exposed to 400B tokens of code, after which it undergoes further training on 150B tokens of math-related data;
    \item \textbf{General Training for 400B Tokens $\rightarrow$ Math Training for 150B Tokens}: For comparison, we also conduct an experiment where the initial 400B tokens are sampled from a broad general corpus curated by DeepSeek-AI, replacing code tokens, to analyze whether code-specific pretraining confers a distinct advantage for mathematical reasoning over general-domain data.
\end{itemize}

\textbf{One-Stage Training}

\begin{itemize}[topsep=0pt]
    \item \textbf{Math Training for 150B Tokens}: DeepSeek-LLM 1.3B is trained solely on 150B tokens of mathematical content;
    \item \textbf{Training on a mixture of 400B Code Tokens and 150B Math Tokens}: Noting that performing math training after code pretraining negatively impacts code-related capabilities, we explore the effects of jointly presenting code and math tokens during a single-stage training process to determine whether this approach not only fosters mathematical reasoning but also prevents catastrophic forgetting of coding skills.
\end{itemize}

\paragraph{Results}
Tables \ref{tab:code-math} and \ref{tab:code-math-general-eval} present the effects of the various training strategies on downstream tasks.

Pretraining with code offers substantial benefits for program-assisted mathematical problem-solving, both in the two-stage and one-stage paradigms. In Table \ref{tab:code-math}, two-stage training beginning with code pretraining alone produces marked improvements on tasks such as GSM8K and MATH, when solved via Python programs. Additional math-specific training in the subsequent stage leads to further progress. Notably, the joint (one-stage) code and math training regime helps to avert the catastrophic forgetting observed in sequential approaches, and enhances both coding (Table \ref{tab:code-math-general-eval}) and program-guided mathematical reasoning (Table \ref{tab:code-math}).

Furthermore, code pretraining also elevates mathematical reasoning even without recourse to external tools. When adopting the two-stage approach, preliminary training on code brings measurable gains, and further math-specific training builds upon this foundation to achieve optimal performance. Conversely, merging code and math data into a single training stage appears to diminish mathematical reasoning proficiency without tool use. This may be attributable to the comparatively modest capacity of DeepSeek-LLM 1.3B, which might be insufficient to master both code and mathematics simultaneously within a joint training regime.

\subsubsection{ArXiv Papers Show Limited Efficacy in Enhancing Mathematical Reasoning}

ArXiv papers are frequently incorporated into the pre-training datasets for mathematical models \citep{minerva,gpt-f,llemma,mathpile}. Despite their widespread usage, there has been little thorough investigation into how these papers affect the development of mathematical reasoning abilities. Surprisingly, our results suggest that arXiv papers are not particularly beneficial for enhancing mathematical reasoning skills. We assess this effect by training models of varying capacities, including DeepSeek-LLM 1.3B and DeepSeek-Coder-Base-v1.5 7B \citep{deepseek-coder}, on arXiv-based datasets subjected to different cleaning and processing protocols:

\begin{itemize}[topsep=0pt]
    \item \textbf{MathPile} \citep{mathpile}: an 8.9B-token collection assembled using heuristic cleaning and filtering processes, of which more than 85\% originates from scientific arXiv manuscripts;
    \item \textbf{ArXiv-RedPajama} \citep{redpajama}: the full set of arXiv LaTeX submissions, with preambles, comments, macros, and bibliography entries omitted, comprising 28.0B tokens.
\end{itemize}

For our analysis, we train DeepSeek-LLM 1.3B on each arXiv dataset for 150B tokens and DeepSeek-Coder-Base-v1.5 7B for 40B tokens. The results indicate that exposure to arXiv content fails to yield improvements in mathematical reasoning and, in some instances, may even impair performance across a range of mathematical benchmarks. These include quantitative problem-solving tasks such as GSM8K and MATH (refer to Table \ref{tab:arxiv-cot}), multiple-choice assessments including MMLU-STEM (Table \ref{tab:arxiv-cot}), and formal mathematics tests like miniF2F (see Table \ref{tab:arxiv_atp}).

Nevertheless, several important caveats temper these findings. Notably, our study does not explore:

\begin{itemize}[topsep=0pt]
    \item The influence of arXiv-derived data on specific mathematical tasks outside our benchmark scope, such as transforming formal proofs or statements into more informal versions (the informalization of theorems);
    \item The interaction effects when arXiv documents are combined with alternative data sources during training;
    \item Whether potential advantages of arXiv material might emerge at larger model scales.
\end{itemize}

Consequently, more comprehensive investigation is necessary, which we leave to subsequent work.

\subsection{Insights of Reinforcement Learning}
\subsubsection{Towards a Unified Paradigm} This section introduces a unified analytical framework to understand various training algorithms—specifically, SFT, RFT, DPO, PPO, and GRPO—and presents experimental results probing the core factors of this framework. In general, the parameter gradient $\theta$ for a given training algorithm can be formally expressed as:

\begin{equation}
    \nabla_{\theta}\mathcal{J}_{\textcolor{red}{\mathcal{A}}}(\theta) = \mathbb{E}[\underbrace{(q,o) \sim \textcolor{red}{\mathcal{D}}}_{Data \ Source}]\left( \frac{1}{|o|} \sum_{t=1}^{|o|}  \underbrace{GC_{{\mathcal{A}}}(q, o, t, \textcolor{red}{\pi_{{rf}}})}_{Gradient \ Coefficient}  \nabla_{\theta}\log \pi_{\theta}(o_t | q, o_{<t})\right).
\label{eq:objective}
\end{equation}

This formulation highlights three fundamental elements: 1) \textit{Data Source $\mathcal{D}$}, specifying the dataset utilized for model training; 2) \textit{Reward Function $\pi_{{rf}}$}, which provides the feedback signal guiding learning; and 3) \textit{Algorithm $\mathcal{A}$}, which ingests both the training examples and the associated rewards to produce the gradient coefficient $GC$, dictating the direction and intensity of parameter updates. The following are several prominent techniques, all of which can be mapped onto this unified framework:

\begin{itemize}
    \item  \textbf{Supervised Fine-tuning (SFT)}: In SFT, the pretrained model undergoes additional training on a dataset curated through human annotation.
    \item \textbf{Rejection Sampling Fine-tuning (RFT)}: Here, the SFT model is further fine-tuned using samples generated by itself in response to SFT prompts, where only those responses deemed correct are retained for continued training.
    \item \textbf{Direct Preference Optimization (DPO)}: DPO iteratively improves the SFT model by exposing it to enhanced outputs sampled from the original model, applying a pairwise loss function rooted in preference comparison.
    \item \textbf{Online Rejection Sampling Fine-tuning (Online RFT)}: Contrasting with standard RFT, Online RFT continually updates the policy model initialized with SFT, leveraging new outputs sampled from the evolving policy in real time.
    \item \textbf{PPO/GRPO}: For PPO/GRPO, the policy is seeded with the SFT weights, then updated based on feedback from outputs generated by the live policy model itself.
\end{itemize}

A summary of the elements comprising these methodologies is presented in Table \ref{tab:data_policy}. For an expanded derivation, consult Appendix \ref{app:analysis-rl}.

\paragraph{Observation about Data Source} We categorize the origins of data into online and offline sampling. Online sampling refers to collecting training data generated by the current policy model's interactions during training, whereas offline sampling draws training data from outputs produced by the pre-trained SFT model. RFT and DPO operate under the offline sampling regime, whereas Online RFT and GRPO utilize online sampling.

As depicted in Figure \ref{fig:rl-analysis}, our analysis demonstrates that Online RFT notably surpasses RFT across both evaluated benchmarks. Initially, both methods show similar performance, as Online RFT and RFT models, alongside their respective SFT counterparts, are closely aligned and the data collected reflects only minimal divergences. As training progresses, however, actor-generated data becomes increasingly distinct from the SFT distribution, providing a greater benefit to real-time sampling strategies. Consequently, Online RFT establishes a clear lead in later training stages, highlighting the efficacy of online training methodologies.

\paragraph{Observation about Gradient Coefficient} In these algorithms, the transformation of reward signals into gradient coefficients underpins the parameter updates. Our experiments distinguish reward functions into ‘Rule’ and ‘Model’ categories. The ‘Rule’ approach evaluates responses based on factual correctness, while the ‘Model’ category involves training a reward model (using data labeled by the rule criterion) to score individual outputs. Equations \ref{eq:GC-RFT} and \ref{eq:GC-GRPO} reveal a central distinction between GRPO and Online RFT: GRPO specifically modifies the gradient coefficient as a function of the reward model’s value, enabling tailored reinforcement or suppression in proportion to reward strength. Conversely, Online RFT lacks this nuanced mechanism, refraining from penalizing incorrect responses and consistently applying an equal positive update to all correct answers, irrespective of their individual quality.

As illustrated in Figure \ref{fig:rl-analysis}, GRPO achieves better results than online RFT, underscoring the advantage of modifying positive and negative gradient coefficients. Moreover, GRPO+PS outperforms GRPO+OS, which suggests that employing granular, step-sensitive gradient coefficients yields additional benefits. Beyond this, our analysis extends to iterative RL: we conduct two consecutive iterations during experimentation. Figure \ref{fig:iter-rl} reveals that this iterative process substantially enhances performance, with the greatest improvement observed after the initial iteration.

\subsubsection{Why RL Works?} This study implements reinforcement learning using a subset of instruction tuning data, which leads to substantial gains over the instruction-tuned baseline. To further clarify the effectiveness of reinforcement learning, we compare the Pass@K and Maj@K accuracies for both Instruct and RL-enhanced models across two benchmark datasets. As depicted in Figure \ref{fig:maj-pass}, RL training raises Maj@K accuracy without affecting Pass@K. These observations suggest that reinforcement learning strengthens the overall response distribution, effectively \textbf{raising the likelihood that a correct answer is found among the TopK outputs, rather than fundamentally advancing baseline abilities}. Analogously, prior work such as \citep{wang2023making} has reported a \textbf{misalignment problem} for reasoning within SFT models, and demonstrated that various preference alignment frameworks \citep{yuan2023rrhf,song2023preference,wang2023making} can enhance the reasoning proficiency of these models.

\subsubsection{How to Achieve More Effective RL?}
\label{sec:effec_rl}
We provide empirical evidence that RL is highly effective for mathematical reasoning problems. Our work also proposes a unified framework for interpreting diverse representative training strategies, where each is viewed as either a form of direct or simplified RL. Equation \ref{eq:objective} synthesizes the essential components: Data Source, Algorithm, and Reward Function. Below, we discuss promising avenues for advancing each of these elements.

\paragraph{Data Source} The selection of data source forms the foundational layer for all training paradigms. In the context of RL, this refers to the set of unlabeled questions and the sampled responses generated by the policy model. For this study, we restricted our approach to questions from the instruction tuning phase and relied on simple nucleus sampling to obtain outputs, which may explain why our RL procedure primarily improves Maj@K scores. Future work will investigate the impact of deploying our RL method on out-of-distribution queries, and of utilizing \textbf{more sophisticated sampling (decoding) approaches}, including tree-search-based techniques \citep{yao2023tree}. Additionally, innovations in \textbf{efficient inference methods} \citep{xia-etal-2023-speculative,leviathan2023fast,kwon2023efficient,xia2024unlocking}—which critically determine the effectiveness of policy model exploration—are poised to play a central role.

\paragraph{Algorithms} Algorithms utilize both the data and the associated reward signals to determine the appropriate gradient coefficients, thereby guiding the update of model parameters. Referring to Equation \ref{eq:objective}, current methodologies fundamentally \textbf{TRUST} the feedback provided by the reward function to adjust the conditional likelihood of specific tokens either upwards or downwards. Nevertheless, there is no guarantee that the reward signal maintains high fidelity in all circumstances, particularly for highly intricate tasks. For instance, despite meticulous annotation by trained professionals in the PRM800K datasets \citep{lightman2023let}, roughly 20\% of labels are still inaccurate\footnote{\url{https://github.com/openai/prm800k/issues/12\#issuecomment-1728491852}}. Therefore, our aim is to investigate reinforcement learning algorithms with inherent robustness to noise in reward signals. We posit that approaches focused on \textbf{WEAK-TO-STRONG} \citep{burns2023weak} alignment have the potential to fundamentally reshape the nature of learning algorithms.

\paragraph{Reward Function} The reward function serves as the foundational source of training feedback. Within RL, it is typically implemented as a neural reward model. We identify three pressing avenues for advancing reward models: 1) \textbf{Improving the generalization capability of reward models.} It is essential that reward models can generalize effectively to unforeseen types of questions and to more sophisticated decoding outputs; failure to do so may result in reinforcement learning simply consolidating the distribution produced by large language models, rather than truly advancing their core competencies; 2) \textbf{Capturing the uncertainty of reward models.} Explicitly modeling uncertainty may serve as an important conduit between suboptimal reward models and the weak-to-strong alignment framework; 3) \textbf{Developing high-quality process reward models efficiently} that can deliver detailed, process-level signals during training for tasks requiring complex reasoning \citep{lightman2023let,wang2023math}.

\section{Conclusion, Limitation, and Future Work} 

We introduce DeepSeekMath, which surpasses all open-source competitors on the competition-level MATH benchmark and approaches the proficiency demonstrated by proprietary models. DeepSeekMath~builds upon DeepSeek-Coder-v1.5 7B and is continually trained with 500B tokens, including a substantial 120B math-specific tokens extracted from Common Crawl. Our thorough ablation analyses reveal that web data represent a rich source of high-quality mathematical material, whereas arXiv contributions may be less impactful than previously assumed. Additionally, we propose Group Relative Policy Optimization (GRPO), an adaptation of Proximal Policy Optimization (PPO), which enables notable improvements in mathematical reasoning while requiring less memory. Empirical results confirm that GRPO remains beneficial even after DeepSeekMath-Instruct 7B achieves top benchmark scores. We further propose a unified framework for interpreting various methods and highlight several promising avenues for advancing reinforcement learning.

Despite its strong quantitative reasoning results, DeepSeekMath~is less adept at geometry and theorem-proving tasks compared to closed-source models. For example, our pilot evaluations show difficulties with questions involving triangles and ellipses, suggesting possible pre-training and fine-tuning data selection biases. Moreover, due to its smaller size, DeepSeekMath~does not match GPT-4's performance in few-shot scenarios—while GPT-4 exhibits gains from few-shot input, DeepSeekMath~shows little difference between zero-shot and few-shot settings. Looking ahead, we plan to refine our engineered data selection processes to further enhance the quality of pre-training data. We also intend to pursue research in the directions summarized in Section \ref{sec:effec_rl} to enable more efficient reinforcement learning for LLMs.

\end{document}